%% Cover letter for the Knowledge-Based Systems submission of
%% "HMRS: Hierarchical document management for evidence-traced retrieval".
%% Scope follows Elsevier's cover-letter guidance: under one page; aim and
%% main findings; fit with the journal's aims and scope; novelty and wider
%% implications. Funding, author declarations and reviewer suggestions are
%% entered in Editorial Manager and are not repeated here.
%% AUTHOR ACTION: replace the bracketed placeholders (name, affiliation,
%% address, email, ORCID) before upload. The date is the compile date.
%% A line that starts with "[" must follow "\\{}" (else LaTeX reads it as
%% the optional argument of \\).
\documentclass[11pt,a4paper]{article}
\usepackage[margin=1in]{geometry}
\usepackage[T1]{fontenc}
\usepackage[utf8]{inputenc}
\usepackage{newtxtext}
\usepackage{enumitem}
\usepackage[hidelinks]{hyperref}
\setlength{\parindent}{0pt}
\setlength{\parskip}{3pt}
\pagestyle{empty}

\begin{document}

\textbf{[Corresponding author name]}\\
{}[Department, Institution]\\
{}[Full postal address]\\
Email: [email]\quad ORCID: [ORCID]

\today

The Editors-in-Chief\\
Knowledge-Based Systems

Dear Editors,

\textbf{Re: ``HMRS: Hierarchical document management for evidence-traced retrieval''}

On behalf of all authors, I submit this manuscript for consideration as a
regular paper in Knowledge-Based Systems. The work is original, has not
been published elsewhere and is not under consideration by another
journal; all authors have approved this submission.

Retrieval from evolving document collections depends on where documents
are placed, how they are described and whether their indexes are current,
yet acquisition, organisation and retrieval are usually studied apart.
HMRS, the Hierarchical Management Retrieval System, is a self-hosted
platform that connects them: it manages collections, documents and
segments, exposes them to agents as catalogue, shelf, document and
segment views through the Model Context Protocol, and adds an optional
server-side evidence gate and post-operation retrieval probes.

The paper's main findings are as follows.
\begin{itemize}[leftmargin=1.5em,itemsep=1pt,topsep=2pt,parsep=0pt]
\item In an early test on a disposable collection, a rename left stale
vector entries. Forced reindexing restored the expected count.
\item The description-bearing route-card configuration selected the labelled
collection on 72/72 questions in all three runs. The nominally name-only
configuration selected it on 37/72 in the first run and 49/72 in each rerun.
\item A configuration combining scored document text with
vector-retrieved passages contained the annotated target string in 52/72
evidence packs, against 37/72 for the reference configuration.
\item In a replay of 16 logged candidate pools from one corpus, fusion
ordering kept a gold document in the top five in 15 pools, against seven
under verifier-score ordering.
\item In the unified blind pass, a judge from a different model family
rated all ten answers from the deployed HMRS agent correct. None of the
paired contrasts against the three baselines built on the same corpus was
significant.
\end{itemize}

Scope conditions are reported with each comparison: the blind grid is
exploratory, reading budgets and indexed corpora differed between arms,
the agent searched outside its five configured collections in nine of ten
questions. The native comparison does not isolate retrieval-backend
performance.

The paper suits the journal's aim of encouraging
``knowledge-based intelligence models, methods, systems, and software
tools'', and continues recent work published in the journal on retrieval
for enterprise collections. Its novelty lies not in a new ranking
algorithm but in treating operator-maintained collection state as a
first-class, inspectable input to agent retrieval; three findings carry
beyond this system: probe the query path after lifecycle writes, evaluate
evidence admission and ordering separately, and report search coverage
alongside results.

Code and data are available at \url{https://github.com/kingdol666/rag-knowledge}.

The submission includes the \LaTeX{} source, highlights and supplementary
material. Thank you for considering our manuscript.

Yours sincerely,

\medskip
[Signature]\\
\textbf{[Name]}, on behalf of all authors\\
{}[Affiliation]\quad [Email]

\end{document}
